\documentclass[12pt]{article}
\usepackage{graphicx} % Required for inserting images
\usepackage[margin=0.5in]{geometry}
\usepackage{amsmath, amssymb}
\usepackage{mathrsfs}


\usepackage{soul}


% Dr. David Brown Commands
\newcommand{\ds}{\displaystyle}
\newcommand{\R}{\mathbb{R}}
\newcommand{\N}{\mathbb{N}}
\newcommand{\Q}{\mathbb{Q}}
\newcommand{\C}{\mathbb{C}}


% Commands to get script r
\usepackage{calligra}
\DeclareMathAlphabet{\mathcalligra}{T1}{calligra}{m}{n}
\DeclareFontShape{T1}{calligra}{m}{n}{<->s*[2.2]callig15}{}
\newcommand{\scripty}[1]{\ensuremath{\mathcalligra{#1}}}

% Unit commands
\newcommand{\degree}{\text{°}}
\newcommand{\unitSpeed}{\frac{\text{m}}{\text{s}}}
\newcommand{\unitAcceleration}{\frac{\text{m}}{\text{s}^2}}

% Constant commands
\newcommand{\avogadro}{6.022\times10^{23}}

\newcommand{\boltzmannConstK}{1.381\times10^{-23}\frac{\text{J}}{\text{K}}}
\newcommand{\boltzmannConstInEV}{8.617\times10^{-5}\frac{\text{eV}}{\text{K}}}
\newcommand{\planckConst}{6.626\times10^{-34}\text{J}\cdot\text{s}}

\newcommand{\atmP}{1.01325\times10^5\,\text{Pa}}

% Known Value Commands
\newcommand{\electronMass}{9.109\times10^{-31}\,\text{kg}}
\newcommand{\protonMass}{1.673\times10^{-27}\,\text{kg}}

% Commands to easily write partial derivatives
\newcommand{\partDeriv}[2]{\frac{\partial #1}{\partial #2}}
\newcommand{\doublePartDeriv}[2]{\frac{\partial^2 #1}{\partial^2 #2}}


\title{Problem 2.3}
\author{Gabriel Decker}


\begin{document}


\maketitle
\begin{flushleft}

In this problem, we will find the electric field a distance $z$ above one end of a line charge of length $L$ and charge density $\lambda$, according to Figure 2.7.
We must also check that the answer makes sense if $z\gg L$.\\~\\

First, we'll define our coordiante system such that the end of the line charge that we're finding the $\vec{E}$ above is at the origin, and the line extends into the $\hat{x}$ direction and our point of interest is in the $\hat{z}$ direction.
This implies that our line runs from $\vec{r'}=0\hat{x}+0\hat{z}$ to $\vec{r'}=L\hat{x}+0\hat{z}$, and our point of interest has coordinates $\vec{r}=0\hat{x}+z\hat{z}$.
Since we're dealing with a line charge, we will use the following equation of an electric field.
$$\vec{E}(\vec{r})=\frac{1}{4\pi\epsilon_0}\int\frac{\lambda(\vec{r'})}{\scripty{r}^2}\hat{\scripty{r}}\,d\ell'$$\\~\\

We will parameterize our line segments by simply representing the different line segments $d\ell'$ as $dx$ and run along the varying $x$ coordinates from $x=0$ to $x=L$.
Also, our line charge density is constant across our line charge, so the following is our integral.
$$\implies\vec{E}(\vec{r})=\frac{1}{4\pi\epsilon_0}\lambda\int_0^L\frac{\hat{\scripty{r}}}{\scripty{r}^2}\,dx$$
$$\implies\vec{E}(\vec{r})=\frac{1}{4\pi\epsilon_0}\lambda\int_0^L\frac{\vec{\scripty{r}}}{\scripty{r}^3}\,dx$$\\~\\

For a given line segment $dx$, it has position $\vec{r'}=x\hat{x}+0\hat{z}$, so $\vec{\scripty{r}}=\vec{r}-\vec{r'}=-x\hat{x}+z\hat{z}$ and $\scripty{r}=\sqrt{x^2+z^2}$.
Plugging this in, we get the following.
$$\implies\vec{E}(\vec{r})=\frac{1}{4\pi\epsilon_0}\lambda\int_0^L\frac{-x\hat{x}+z\hat{z}}{(x^2+z^2)^{3/2}}\,dx$$
$$\implies\vec{E}(\vec{r})=\frac{1}{4\pi\epsilon_0}\lambda\left[-\hat{x}\int_0^L\frac{x}{(x^2+z^2)^{3/2}}\,dx+z\hat{z}\int_0^L\frac{1}{(x^2+z^2)^{3/2}}\,dx\right]$$
We now have to solve the two inner integrals.\\~\\

$$\int_0^L\frac{x}{(x^2+z^2)^{3/2}}\,dx$$
Using the substitution $u=x^2+z^2\implies du=2x\,dx$.
$$\implies\int_0^L\frac{x}{(x^2+z^2)^{3/2}}\,dx=\frac{1}{2}\int_{z^2}^{L^2+z^2}\frac{1}{u^{3/2}}\,du$$
$$\implies\int_0^L\frac{x}{(x^2+z^2)^{3/2}}\,dx=\frac{1}{2}\int_{z^2}^{L^2+z^2}u^{-3/2}\,du$$
$$\implies\int_0^L\frac{x}{(x^2+z^2)^{3/2}}\,dx=-u^{-1/2}\Big|_{z^2}^{L^2+z^2}$$
$$\implies\int_0^L\frac{x}{(x^2+z^2)^{3/2}}\,dx=-\frac{1}{\sqrt{u}}\Big|_{z^2}^{L^2+z^2}$$
$$\implies\int_0^L\frac{x}{(x^2+z^2)^{3/2}}\,dx=-\left(\frac{1}{\sqrt{L^2+z^2}}-\frac{1}{\sqrt{z^2}}\right)$$\\~\\

Now for the other integral.
$$\int_0^L\frac{1}{(x^2+z^2)^{3/2}}\,dx$$
For this, there is the following solution that comes from doing a trig substitution $x=z\tan(\theta)$.
$$\int\frac{1}{(x^2+z^2)^{3/2}}\,dx=\frac{x}{z^2\sqrt{x^2+z^2}}$$
So,
$$\implies\int_0^L\frac{1}{(x^2+z^2)^{3/2}}\,dx=\frac{L}{z^2\sqrt{L^2+z^2}}$$\\~\\

Plugging in our results, we get the following electric field.
$$\vec{E}(\vec{r})=\frac{1}{4\pi\epsilon_0}\lambda\left[-\hat{x}\int_0^L\frac{x}{(x^2+z^2)^{3/2}}\,dx+z\hat{z}\int_0^L\frac{1}{(x^2+z^2)^{3/2}}\,dx\right]$$
$$\implies\vec{E}(\vec{r})=\frac{1}{4\pi\epsilon_0}\lambda\left[\left(\frac{1}{\sqrt{L^2+z^2}}-\frac{1}{\sqrt{z^2}}\right)\hat{x}+\frac{L}{z\sqrt{L^2+z^2}}\hat{z}\right]$$\\~\\

If $z\gg L$, then $\sqrt{L^2+z^2}=\sqrt{z^2(\frac{L^2}{z^2}+1)}\approx\sqrt{z^2}=z$.
This affects $\vec{E}$ in the following way.
$$\implies\vec{E}(\vec{r})=\frac{1}{4\pi\epsilon_0}\lambda\left[\left(\frac{1}{z}-\frac{1}{z}\right)\hat{x}+\frac{L}{z^2}\hat{z}\right]$$
$$\implies\vec{E}(\vec{r})=\frac{1}{4\pi\epsilon_0}\frac{\lambda L}{z^2}\hat{z}$$
Since $\lambda$ is the charge density, $\lambda L$ is the total charge of the line, so the above equation is functionally the same as if the line segment were a point charge with charge $\lambda L$.


\end{flushleft}

\end{document}
